% GENERATED FILE. Edit canonical lessons, then run npm run generate:books.
% canonical-source-hash: fnv1a64:618d6b5e20f6d3cd
% canonical-lessons: MW-C204-ghumno, MW-C204-bhir, MW-C204-maidan, MW-C204-kinaro, MW-C204-nakso

\chapter{Out and About}
\label{ch:mw-out-and-about}
\hlchaptermodality{204}

\begin{chapteropening}
\textbf{By the end of this chapter:} \emph{I can stroll round the town and find my way.}

Complete the last lesson of chapter 204: I can stroll round the town and find my way.
\end{chapteropening}

\section[\texorpdfstring{\mw{घूमणो} (ghūmṇo)}{घूमणो (ghūmṇo)}]{\mw{घूमणो} (ghūmṇo) — to stroll; to go round}

\label{lesson:MW-C204-ghumno}



\begin{quote}
Before the new one: say the Marwadi for to lose one's way, then the Marwadi for to be visible.
\end{quote}

\subsection*{You'll want to know: \mw{घूमणो}}
\textbf{\mw{घूमणो}} — \emph{ghūmṇo} — \textquotedblleft{}to stroll; to go round\textquotedblright{}.

\begin{quote}
\mw{आथण} \mw{हूं} \mw{बाजार} \mw{मांय} \mw{घूमूं} \mw{हूं।}
\end{quote}

\begin{quote}
\emph{āthaṇ hū̃ bājār mā̃y ghūmū̃ hū̃.}
\end{quote}

\begin{quote}
In the evening I stroll round the market.
\end{quote}

Past (I …): \mw{हूं} \mw{सहर} \mw{मांय} \mw{घूम्यो}, \emph{hū̃ sahar mā̃y ghūmyo}.

\begin{cousinweb}
\textbf{\mw{घूमणो}} (\emph{ghūmṇo}) is inherited from a reconstructed \emph{*ghummati}, \textquotedblleft{}revolves\textquotedblright{}. CDIAL gives no Marwari form in this entry, but lists Hindi \emph{ghūmnā}, \textquotedblleft{}to revolve, wander\textquotedblright{} (CDIAL 4485).
\end{cousinweb}

\begin{grammarlens}[title={What you've built}]
One more word for this chapter.
\end{grammarlens}

\subsection*{Your turn}
\begin{itemize}
\raggedright
  \item \emph{Say it:} \emph{ghūmṇo}
  \item \emph{Say it:} \emph{ghūmṇo}, once more
  \item \emph{Say it:} \emph{āthaṇ hū̃ bājār mā̃y ghūmū̃ hū̃.} — In the evening I stroll round the market.
  \item \emph{Recall:} say the Marwadi for to lose one's way, then the Marwadi for to be visible, then say \emph{ghūmṇo} again
\end{itemize}

\begin{tcolorbox}[breakable,colback=teal!4,colframe=teal!35!black,title={Before you move on}]
What does \mw{घूमणो} mean? (\textquotedblleft{}To stroll; to go round\textquotedblright{}.) Say it once more.
\end{tcolorbox}
\section[\texorpdfstring{\mw{भीड़} (bhīṛ)}{भीड़ (bhīṛ)}]{\mw{भीड़} (bhīṛ) — a crowd}

\label{lesson:MW-C204-bhir}



\begin{quote}
Before the new one: say the Marwadi for to be visible, then the Marwadi for to stroll.
\end{quote}

\subsection*{You'll want to know: \mw{भीड़}}
\textbf{\mw{भीड़}} — \emph{bhīṛ} — \textquotedblleft{}a crowd\textquotedblright{}.

\begin{quote}
\mw{आज} \mw{बाजार} \mw{मांय} \mw{घणी} \mw{भीड़} \mw{है।}
\end{quote}

\begin{quote}
\emph{āj bājār mā̃y ghaṇī bhīṛ hai.}
\end{quote}

\begin{quote}
The market is very crowded today.
\end{quote}

A feminine noun.

\begin{cousinweb}
\textbf{\mw{भीड़}} (\emph{bhīṛ}) is inherited from a reconstructed \emph{*bhiṭ}, \textquotedblleft{}meet, throng\textquotedblright{}; CDIAL lists Old Marwari \emph{bheṭi}, \textquotedblleft{}meeting\textquotedblright{}, and Hindi \emph{bhīṛ}, \textquotedblleft{}crowd\textquotedblright{}, under it (CDIAL 9490).
\end{cousinweb}

\begin{grammarlens}[title={What you've built}]
One more word for this chapter.
\end{grammarlens}

\subsection*{Your turn}
\begin{itemize}
\raggedright
  \item \emph{Say it:} \emph{bhīṛ}
  \item \emph{Say it:} \emph{bhīṛ}, once more
  \item \emph{Say it:} \emph{āj bājār mā̃y ghaṇī bhīṛ hai.} — The market is very crowded today.
  \item \emph{Recall:} say the Marwadi for to be visible, then the Marwadi for to stroll, then say \emph{bhīṛ} again
\end{itemize}

\begin{tcolorbox}[breakable,colback=teal!4,colframe=teal!35!black,title={Before you move on}]
What does \mw{भीड़} mean? (\textquotedblleft{}A crowd\textquotedblright{}.) Say it once more.
\end{tcolorbox}
\section[\texorpdfstring{\mw{मैदान} (maidān)}{मैदान (maidān)}]{\mw{मैदान} (maidān) — an open ground}

\label{lesson:MW-C204-maidan}



\begin{quote}
Before the new one: say the Marwadi for to stroll, then the Marwadi for a crowd.
\end{quote}

\subsection*{You'll want to know: \mw{मैदान}}
\textbf{\mw{मैदान}} — \emph{maidān} — \textquotedblleft{}an open ground\textquotedblright{}.

\begin{quote}
\mw{मैदान} \mw{स्कूल} \mw{कनै} \mw{है।}
\end{quote}

\begin{quote}
\emph{maidān skūl kanai hai.}
\end{quote}

\begin{quote}
The open ground is near the school.
\end{quote}

A masculine noun.

\begin{cousinweb}
\textbf{\mw{मैदान}} (\emph{maidān}) is usually traced to Persian \emph{maidān}, \textquotedblleft{}an open plain\textquotedblright{}.
\end{cousinweb}

\begin{grammarlens}[title={What you've built}]
One more word for this chapter.
\end{grammarlens}

\subsection*{Your turn}
\begin{itemize}
\raggedright
  \item \emph{Say it:} \emph{maidān}
  \item \emph{Say it:} \emph{maidān}, once more
  \item \emph{Say it:} \emph{maidān skūl kanai hai.} — The open ground is near the school.
  \item \emph{Recall:} say the Marwadi for to stroll, then the Marwadi for a crowd, then say \emph{maidān} again
\end{itemize}

\begin{tcolorbox}[breakable,colback=teal!4,colframe=teal!35!black,title={Before you move on}]
What does \mw{मैदान} mean? (\textquotedblleft{}An open ground\textquotedblright{}.) Say it once more.
\end{tcolorbox}
\section[\texorpdfstring{\mw{किनारो} (kināro)}{किनारो (kināro)}]{\mw{किनारो} (kināro) — an edge; a shore}

\label{lesson:MW-C204-kinaro}



\begin{quote}
Before the new one: say the Marwadi for a crowd, then the Marwadi for an open ground.
\end{quote}

\subsection*{You'll want to know: \mw{किनारो}}
\textbf{\mw{किनारो}} — \emph{kināro} — \textquotedblleft{}an edge; a shore\textquotedblright{}.

\begin{quote}
\mw{किनारो} \mw{अठै} \mw{सूं} \mw{दूर} \mw{कोनी।}
\end{quote}

\begin{quote}
\emph{kināro aṭhai sū̃ dūr konī.}
\end{quote}

\begin{quote}
The shore isn't far from here.
\end{quote}

A masculine noun.

\begin{cousinweb}
\textbf{\mw{किनारो}} (\emph{kināro}) is usually traced to Persian \emph{kinār}, \textquotedblleft{}edge, shore\textquotedblright{}.
\end{cousinweb}

\begin{grammarlens}[title={What you've built}]
One more word for this chapter.
\end{grammarlens}

\subsection*{Your turn}
\begin{itemize}
\raggedright
  \item \emph{Say it:} \emph{kināro}
  \item \emph{Say it:} \emph{kināro}, once more
  \item \emph{Say it:} \emph{kināro aṭhai sū̃ dūr konī.} — The shore isn't far from here.
  \item \emph{Recall:} say the Marwadi for a crowd, then the Marwadi for an open ground, then say \emph{kināro} again
\end{itemize}

\begin{tcolorbox}[breakable,colback=teal!4,colframe=teal!35!black,title={Before you move on}]
What does \mw{किनारो} mean? (\textquotedblleft{}An edge; a shore\textquotedblright{}.) Say it once more.
\end{tcolorbox}
\section[\texorpdfstring{\mw{नक्सो} (nakso)}{नक्सो (nakso)}]{\mw{नक्सो} (nakso) — a map}

\label{lesson:MW-C204-nakso}



\begin{quote}
Before the new one: say the Marwadi for an open ground, then the Marwadi for an edge.
\end{quote}

\subsection*{You'll want to know: \mw{नक्सो}}
\textbf{\mw{नक्सो}} — \emph{nakso} — \textquotedblleft{}a map\textquotedblright{}.

\begin{quote}
\mw{नक्सो} \mw{जेब} \mw{मांय} \mw{है।}
\end{quote}

\begin{quote}
\emph{nakso jeb mā̃y hai.}
\end{quote}

\begin{quote}
The map is in my pocket.
\end{quote}

A masculine noun.

\begin{cousinweb}
\textbf{\mw{नक्सो}} (\emph{nakso}) is usually traced, through Persian, to Arabic \emph{naqsh}, \textquotedblleft{}a drawing, an engraving\textquotedblright{}.
\end{cousinweb}

\begin{grammarlens}[title={What you've built}]
One more word for this chapter.
\end{grammarlens}

\subsection*{Your turn}
\begin{itemize}
\raggedright
  \item \emph{Say it:} \emph{nakso}
  \item \emph{Say it:} \emph{nakso}, once more
  \item \emph{Say it:} \emph{nakso jeb mā̃y hai.} — The map is in my pocket.
  \item \emph{Recall:} say the Marwadi for an open ground, then the Marwadi for an edge, then say \emph{nakso} again
\end{itemize}

\begin{tcolorbox}[breakable,colback=teal!4,colframe=teal!35!black,title={Before you move on}]
What does \mw{नक्सो} mean? (\textquotedblleft{}A map\textquotedblright{}.) Say it once more.
\end{tcolorbox}
